% =====================================================================
% protocolo.tex
% Responsabilidad única: mostrar en un diagrama el protocolo experimental
% de la Parte II (partición 60/20/20, qué se hace con cada partición y en
% qué punto se congelan las decisiones antes de abrir prueba). Los tamaños
% de partición son macros generadas por el notebook.
% =====================================================================
\begin{figure}[H]
\centering
\begin{tikzpicture}[
    every node/.style={font=\small},
    caja/.style={draw, rounded corners=2pt, align=center, minimum height=1.05cm, inner sep=4pt},
    paso/.style={draw, rounded corners=2pt, align=center, fill=gray!6, text width=3.2cm,
                 minimum height=1.7cm, inner sep=4pt},
    flecha/.style={-{Stealth[length=2.2mm]}, thick},
]
% --- Datos y partición -------------------------------------------------
\node[caja, fill=gray!10, text width=14cm] (datos) at (0,0)
    {\textbf{Datos:}\enspace caso A: \texttt{load\_diabetes}, 442 obs.\\
     caso B: \texttt{load\_breast\_cancer}, 569 obs.\ (1 = maligno)\\[1pt]
     {\scriptsize\color{gray!70!black}\texttt{random\_state=42}, estratificado en B, mismos índices para todos los modelos}};

\node[caja, fill=classblue!14, draw=classblue, text width=3.2cm] (tr) at (-5,-2.3)
    {\textbf{Entrenamiento}\\ 60\,\% \\ A: $\NRegTr$ \quad B: $\NClfTr$};
\node[caja, fill=classgreen!16, draw=classgreen, text width=3.2cm] (val) at (0,-2.3)
    {\textbf{Validación}\\ 20\,\% \\ A: $\NRegVal$ \quad B: $\NClfVal$};
\node[caja, fill=classred!12, draw=classred, text width=3.2cm] (te) at (5,-2.3)
    {\textbf{Prueba}\\ 20\,\% \\ A: $\NRegTest$ \quad B: $\NClfTest$};

\foreach \n in {tr, val, te} \draw[flecha, gray] (datos.south -| \n) -- (\n.north);

% --- Qué se hace con cada partición ------------------------------------
\node[paso, below=0.7cm of tr] (ajuste)
    {\textbf{Ajustar}\\ escaladores, coeficientes y árboles (Ej.~6--9)};
\node[paso, below=0.7cm of val] (decidir)
    {\textbf{Decidir}\\ $C$, profundidad, hojas, umbral $t$ y candidato};
\node[paso, below=0.7cm of te] (evaluar)
    {\textbf{Evaluar una sola vez}\\ métricas e IC bootstrap (Ej.~10)};

\draw[flecha, classblue] (tr) -- (ajuste);
\draw[flecha, classgreen] (val) -- (decidir);
\draw[flecha, classred] (te) -- (evaluar);

\draw[flecha] (ajuste) -- node[above, font=\scriptsize, align=center] {modelos\\ ajustados} (decidir);
\draw[flecha] (decidir) -- node[above, pos=0.3, font=\scriptsize, align=center] {candidatos\\ congelados} (evaluar);

% --- Barrera de congelamiento ------------------------------------------
\coordinate (barrera) at ($(te.west)+(-0.28,0)$);
\draw[classred, very thick, dashed] ($(barrera)+(0,0.95)$) -- ($(barrera |- evaluar.south)+(0,-0.15)$);
\node[classred, font=\scriptsize\bfseries, align=center, fill=white, inner sep=1pt]
    at ($(barrera)+(0,1.25)$) {congelado: sin reentrenar\\ ni cambiar de candidato};

% --- Transformaciones aprendidas solo con entrenamiento -----------------
\draw[flecha, dashed, gray] (ajuste.south) to[out=-30, in=-150] (decidir.south);
\draw[flecha, dashed, gray] (ajuste.south) to[out=-35, in=-145] (evaluar.south);
\node[font=\scriptsize, gray!70!black, align=center] at ($(decidir.south)+(0,-2.05)$)
    {escaladores aprendidos solo en entrenamiento y \emph{aplicados} a validación y prueba};
\end{tikzpicture}
\caption{Protocolo que seguimos en la Parte II: cada partición cumple un
único papel (ajustar, decidir o evaluar), los escaladores se aprenden solo
con entrenamiento y todas las decisiones (hiperparámetros, umbrales y
candidatos) las congelamos antes de consultar la prueba, que usamos una
sola vez en el Ejercicio~10.}
\label{fig:protocolo}
\end{figure}
